\scanpage{72}\mspage{63}

de l'existence d'une telle fonction support, peuvent être \unsure{posés|posées} en \unsure{axiomes|axiome} sans inconvénient pratique.

La seule exception à cette règle concerne les géométries de lieux définis \deleted{\unsure{p}} en termes d'une fonction distance, où la disjonction est décrite en termes d'une distance $>\varepsilon$ \inserted{ou $\geq\varepsilon$}, et ne peut sans doute se ramener à une question de \deleted{\unsure{int|l'int}} ``disjonction ensembliste''. De toutes façons, je n'ai pas regardé ce \unsure{cas} d'assez près pour être assuré que le langage géométrique des magasins que je propose ici, s'y applique également.

\uline{Exemples} 1) \deleted{Dans} Dire que dans \unsure{les|des|le|un} magasin $M$, \deleted{\unsure{les}} $\lvert\unsure{\circ|0}\rvert$ est \unsure{définissable|definissable} par $\overset{\circ}{\ll}$, (p.~30, \unsure{31|71}) signifie : que
\[
 X \longrightarrow \mathrm{Omb}(X) \qquad M \longrightarrow \unsure{\mathcal{P}|P}(M)
\]
est une fonction support \inserted{\uline{fidèle}$^{*}$}. \deleted{\unsure{Mais cette|Mais alors|Mais elle}} \unsure{Sa|Ça} en fait un bon paquet !

% slanted
\margin{$^{*}$\unsure{NB} C'est \uuline{\unsure{tjrs}} une fonction support.}

% archive sheet number: 72
